%\section{PAGINA 4}
\noindent incógnitas de este sistema de ecuaciones. Finalmente se supone que se desea hallar 
$I_2$ en $R_3$. Empezamos en A en el lazo 1 y escribimos la ecuación de tensión: 
\[
+ V_1 - (I_1 + I_2) \times R_2 + V_2 - I_1 \times R_1 = 0
\tag{27-6}
\]
\noindent Sustituyendo los valores de circuito en la ecuación (27-6), tenemos:
\[
100 - 2.200 \times (I_1 + I_2) + 150 - 1.000 \times I_1 = 0
\tag{27-7}
\]
\noindent Simplificando, obtenemos:
\[
3.200I_1 + 2.200I_2 = 250
\tag{27-8}
\]

\noindent Ahora empezamos en el punto D del bucle 2. La ecuación de tensión es:
\[
- I_2R_3 - (I_1 + I_2)R_2 + 150 = 0
\tag{27-9}
\]

\noindent Substituyendo los valores del circuito en la ecuación (27-9):
\[
- 1.200I_2 - 2.200(I_1 + I_2) + 150 = 0
\tag{27-10}
\]

\noindent Simplificando, tenemos:
\[
2.200I_1 + 3.400I_2 = 150
\tag{27-11}
\]

\noindent Con el sistema de dos ecuaciones (27-8) y (27-11) hallamos $I_1$ e $I_2$, obteniendo:
\[
\begin{aligned}
I_1 &= 0,0861 \, \text{A} \\
I_2 &= -0,0115 \, \text{A}
\end{aligned}
\]  

\noindent El signo negativo en $I_2$ significa simplemente que la corriente tiene sentido contrario al supuesto inicialmente. $I_2$ es la corriente que circula por $R_3$.

\subsection{Resolución de circuito por métodos de análisis}

\subsection{Enunciado del problema}
Consideremos el circuito de la figura 27-4 donde se desea hallar la corriente $I_2$ que circula por la resistencia $R_3$.

\subsection{Resolución por leyes de Kirchhoff}
Empezamos en el punto A del bucle 1 y escribimos la ecuación de tensión:
\[
+ V_1 - (I_1 + I_2)R_2 + V_2 - I_1 \times R_1 = 0
\tag{27-6}
\]

Substituyendo los valores del circuito en la ecuación (27-6), tenemos:
\[
100 - 2.200(I_1 + I_2) + 150 - 1.000I_1 = 0
\tag{27-7}
\]

Simplificando, obtenemos:
\[
3.200I_1 + 2.200I_2 = 250
\tag{27-8}
\]

Ahora empezamos en el punto D del bucle 2. La ecuación de tensión es:
\[
- I_2R_3 - (I_1 + I_2)R_2 + 150 = 0
\tag{27-9}
\]

Substituyendo los valores del circuito en la ecuación (27-9):
\[
- 1.200I_2 - 2.200(I_1 + I_2) + 150 = 0
\tag{27-10}
\]

Simplificando, tenemos:
\[
2.200I_1 + 3.400I_2 = 150
\tag{27-11}
\]

Con el sistema de dos ecuaciones (27-8) y (27-11) hallamos $I_1$ e $I_2$, obteniendo:
\[
\begin{aligned}
I_1 &= 0,0861 \, \text{A} \\
I_2 &= -0,0115 \, \text{A}
\end{aligned}
\]  

\noindent El signo negativo en $I_2$ significa simplemente que la corriente tiene sentido contrario al supuesto inicialmente. $I_2$ es la corriente que circula por $R_3$.

\section*{Resolución de circuito por métodos de análisis}

\subsection{Enunciado del problema}
Consideremos el circuito de la figura 27-4 donde se desea hallar la corriente $I_2$ que circula por la resistencia $R_3$.

\subsection{Resolución por leyes de Kirchhoff}
Empezamos en el punto A del bucle 1 y escribimos la ecuación de tensión:
\[
+ V_1 - (I_1 + I_2)R_2 + V_2 - I_1 \times R_1 = 0
\tag{27-6}
\]

Substituyendo los valores del circuito en la ecuación (27-6), tenemos:
\[
100 - 2.200(I_1 + I_2) + 150 - 1.000I_1 = 0
\tag{27-7}
\]

Simplificando, obtenemos:
\[
3.200I_1 + 2.200I_2 = 250
\tag{27-8}
\]

Ahora empezamos en el punto D del bucle 2. La ecuación de tensión es:
\[
- I_2R_3 - (I_1 + I_2)R_2 + 150 = 0
\tag{27-9}
\]

Substituyendo los valores del circuito en la ecuación (27-9):
\[
- 1.200I_2 - 2.200(I_1 + I_2) + 150 = 0
\tag{27-10}
\]

Simplificando, tenemos:
\[
2.200I_1 + 3.400I_2 = 150
\tag{27-11}
\]

Con el sistema de dos ecuaciones (27-8) y (27-11) hallamos $I_1$ e $I_2$, obteniendo:
\[
\begin{aligned}
I_1 &= 0,0861 \, \text{A} \\
I_2 &= -0,0115 \, \text{A}
\end{aligned}
\]

El signo negativo en $I_2$ significa simplemente que la corriente tiene sentido contrario al supuesto inicialmente. $I_2$ es la corriente que circula por $R_3$.

\subsection{Resolución por el teorema de Thévenin}
Para hallar $I_2$ en $R_3$ aplicando el teorema de Thévenin, primero calculamos $V_{TH}$, que es la tensión en circuito abierto entre los puntos F y G.

La tensión de circuito abierto entre F y G es $V_{FG} = V_{TH} = -150 + I_3R_2$. Por tanto, necesitamos hallar $I_3$ antes de determinar $V_{TH}$.

Aplicando la ley de tensión de Kirchhoff al bucle ABCDA:
\[
100 - I_3(2.200) + 150 - I_3(1.000) = 0
\tag{27-12}
\]

Resolviendo:
\[
\begin{aligned}
3.200I_3 &= 250 \\
I_3 &= \frac{250}{3.200} = 0,0781 \, \text{A}
\end{aligned}
\]

Entonces:
\[
I_3R_2 = 0,0781 \times 2.200 = 171,82 \, \text{V}
\]

La tensión Thévenin es:
\[
V_{TH} = -150 + I_3R_2 = 21,82 \, \text{V}
\tag{27-13}
\]

Para hallar $R_{TH}$, substituimos $V_1$ y $V_2$ por sus resistencias internas (supuestas iguales a cero) y hallamos la resistencia entre F y G:
\[
R_{TH} = \frac{2.200 \times 1.000}{2.200 + 1.000} = 687,5 \, \Omega
\tag{27-14}
\]

Finalmente, hallamos $I_2$ en el circuito equivalente de Thévenin:
\[
I = I_2 = \frac{V_{TH}}{R_{TH} + R_3} = \frac{21,82}{687,5 + 1.200} = \frac{21,82}{1.887,5} = 0,0115 \, \text{A}
\tag{27-15}
\]

\subsection{Resolución por el teorema de Norton}
Para hallar $I_2$ en $R_3$ por el teorema de Norton, desarrollamos el circuito equivalente como se muestra en las figuras 27-5b y 27-5c.

En la figura 27-5b, $R_3$ está reemplazada por un cortocircuito. $I_N$, la corriente de cortocircuito entre los terminales F y G, es la corriente del generador equivalente de Norton.

Para hallar $I_N$ aplicamos las leyes de Kirchhoff a los bucles 1 y 2 de la figura 27-5b:
\begin{itemize}
    \item Bucle 1: Trayecto ABCGFDA
    \item Bucle 2: Trayecto DFGCD
\end{itemize}

\subsection{Conclusiones}
\begin{itemize}
    \item Se obtiene el mismo resultado ($I_2 = 0,0115 \, \text{A}$) por los tres métodos
    \item El teorema de Thévenin y Norton proporcionan métodos sistemáticos para simplificar circuitos complejos
    \item La elección del método depende de la configuración específica del circuito y de lo que se necesite calcular
\end{itemize}

\begin{table}[H]
    \centering
    \caption{Comparación de resultados obtenidos por diferentes métodos}
    \label{tab:comparacion_metodos}
    \begin{tabular}{lcc}
        \toprule
        \textbf{Método} & \textbf{Corriente $I_2$ (A)} & \textbf{Observaciones} \\
        \midrule
        Leyes de Kirchhoff & 0,0115 & Método general pero laborioso \\
        Teorema de Thévenin & 0,0115 & Útil para análisis de un elemento \\
        Teorema de Norton & 0,0115 & Equivalente a Thévenin, diferente formulación \\
        \bottomrule
    \end{tabular}
\end{table}

\textbf{Nota:} Los valores numéricos corresponden a resistencias en ohmios, tensiones en voltios y corrientes en amperios.